\documentclass[times, a4paper, 12pt, onecolumn, oneside]{article} 

\usepackage[a4paper,left=2.0cm,right=2.0cm,bottom=1.0cm,top=2.0cm,includefoot,includehead]{geometry}
\usepackage{amsmath}
\usepackage{amsthm}
\usepackage{amssymb}
\usepackage{amsfonts}
\usepackage{graphicx} 
\usepackage{fancyhdr}
\usepackage{natbib}

\pagestyle{fancyplain}
\date{}
\lhead{{\em Immonen (LUT) }}
\chead{Exercises Week 2}
\rhead{\today}

\bibliographystyle{apalike}

\title{ Exercises Week 2}
\author{Veikka Immonen}
\date{\today}

\begin{document}

\maketitle

\section{Tweaking $\nu$ and $\ell$}

\subsection{One realisation of Matèrn process with $\nu=1.75$}

Presented in Figure~\ref{fig:realizations}. A "true"
process is taken using a grid of 1001 points. The
"observed data" contains 51 evenly spaced points taken
from the original process. Moreover, white noise
$\epsilon \sim \mathcal{N}(0, 0.1^2)$ is added to the
observations. As length scale was not specified here, so
it is set to $\ell = 0.1$

\begin{figure}
    \centering
    \includegraphics[width=\textwidth]{./realisation_1d_51_1001.pdf}
    \caption{One realisation of Matèrn process with $\nu=1.75$ and $\ell=0.1$. }
    \label{fig:realizations}
\end{figure}

\subsection{Experimenting with $\nu=1, ...,5$}

Using a grid with a step size of 0.25. The kernel is kept the same otherwise, i.e., $\ell=0.1$.
Results presented in Figure~\ref{fig:nu_fits}. Mispecification in this case refers to
any $\nu \ne p + \frac{1}{2},\ \mathbb{N}_+$. While nothing out of ordinary wasn't
observed during computing the results, such $\nu$ is practice is not the best choice,
due to huge increase in computing complexity in Matèrn covariance.

\begin{figure}
    \centering
    \includegraphics[width=\textwidth]{./nu_fits_51_1001.pdf}
    \caption{GPr fits with different $\nu$ values.}
    \label{fig:nu_fits}
\end{figure}

\subsection{RMSE plot of $\nu$.}

Presented in Figure~\ref{fig:nu_rmse}. Based on the RMSE, the best fit were obtained using $\nu=1.75$.
However, the estimates with other $\nu$ are still satisfactory.

\begin{figure}
    \centering
    \includegraphics[width=\textwidth]{./rmse_nu_1d_51_1001.pdf}
    \caption{RMSE w.r.t. $\nu$.}
    \label{fig:nu_rmse}
\end{figure}

\subsection{Same experiments, but with $\ell$}

Tested on grid with $\ell = 10^{\{-3, -2.75, ..., 1\}}$, and $\nu=1.75$.

All fits are presented in Figure~\ref{fig:ell_fits}, and RMSE plot is presented in Figure~\ref{fig:ell_rmse}.
There is much more contrast between results with different $\ell$ values. A conclusion can be made: finetune $\ell$, not $\nu$.

\begin{figure}
    \centering
    \includegraphics[width=\textwidth]{./ell_fits_51_1001.pdf}
    \caption{GPr fits with different $\ell$ values.}
    \label{fig:ell_fits}
\end{figure}


\begin{figure}
    \centering
    \includegraphics[width=\textwidth]{./rmse_ell_1d_51_1001.pdf}
    \caption{RMSE w.r.t. $\nu$.}
    \label{fig:ell_rmse}
\end{figure}

\section{MCMC on both state and $\nu$}

Now, the hierarchy in the system goes:

$$
\nu \longrightarrow K_\nu \longrightarrow X_* \longrightarrow X
$$

therefore

$$
p(\nu, X_* | X) \propto p(X | X_*, \nu) p(X_* | \nu) p(\nu)
$$

where

$$
p(X | X_*, \nu) = p(X | X_*) = \mathcal{N}(HX_*, \sigma_\epsilon^2 I),
$$

and

$$
p(X_* | \nu) = \mathcal{N}(0, K_\nu).
$$

Since $X_*$ is conditioned to $\nu$, blocking MCMC is performed,
where $X_*$ and $\nu$ are sampled alternately using pCN and MH,
respectively. Moreover, logarithm $z =\log_{10}(\nu)$ is sampled,
to preserve values of $\nu$ in the positive range. The steps are roughly

\begin{enumerate}
\item Draw $X_*' \sim \mathcal{N}\left(\sqrt{1 - \beta_{X_*}^2} X_*^{(k)}, \beta_{X_*}^2 K_\nu\right)$ (pCN)
\item Set $X_*^{(k+1)} = X_*'$ with probability $\min\left\{1, \frac{p(X | X_*')}{p(X | X_*^{(k)})}\right\}$, otherwise to $X_*^{(k)}$
\item Draw $z' \sim \mathcal{N} \left( z^{(k)}, \beta_{\nu}^2 \right)$ (MH)
\item Set $z^{(k + 1)} = z'$ with probability $\min\left\{1, \frac{p(z', X_*^{(k+1)} | X)}{p(z^{(k)}, X_*^{(k+1)} | X)}\right\}$, otherwise to $z^{(k)}$,
\end{enumerate}
assuming a Gaussian (symmetric) random walk in both proposals. Note, that
$$
\frac{p(z', X_*^{(k+1)} | X)}{p(z^{(k)}, X_*^{(k+1)} | X)} = \frac{p(X | X_*^{(k+1)}) p(X_*^{(k+1)} | z') p(z')}{p(X | X_*^{(k+1)}) p(X_*^{(k+1)} | z^{(k)}) p(z^{(k)})} = \frac{p(X_*^{(k+1)} | z') p(z')}{p(X_*^{(k+1)} | z^{(k)}) p(z^{(k)})}
$$
and
\begin{itemize}
\item $p(z) \propto \exp\left( -\frac{z^2}{2\sigma_z^2}\right)$
\item $p(X_* | z) \propto \exp\left(-\frac{1}{2} (X_*^\intercal K_{e^z}^{-1} X_* + \log |K_{e^z}|)\right)$
\item $p(X | X_*) \propto \exp\left(-\frac{1}{2\sigma_\epsilon^2} ||X - HX_*||^2\right)$
\item $\beta_{X_*} = \beta_\nu = 0.01$
\item $\sigma_\nu = 1.0$
\item $\text{Length of the chain} = 10000$
\end{itemize}

Also, dimensions were reduced to $n_*=101$ (true process) and $n=21$ (data).
The sampled distribution of $\nu$ is presented in Figure~\ref{fig:nu_dist},
and the mean posterior fit to the data is presented in Figure~\ref{fig:mcmc}.

\begin{figure}
    \centering
    \includegraphics[width=\textwidth]{./nu_chain.pdf}
    \caption{Chain and histogram of $\nu$.}
    \label{fig:nu_dist}
\end{figure}


\begin{figure}
    \centering
    \includegraphics[width=\textwidth]{./mcmc.pdf}
    \caption{MCMC fit of the state.}
    \label{fig:mcmc}
\end{figure}


% \bibliography{../template/references}

\end{document}
